%%%%%%%%%%%%%%%%%%%%%%%%%%%%%%%%%%%%%%%%%%%%%%%%%%%%%%%%%%%%%%%%%%%%%%%%
% Uni Duesseldorf
% Lehrstuhl fuer Datenbanken und Informationssysteme
% Vorlage fuer Bachelor-/Masterarbeiten
% Optimiert fuer den Original-Latex-Kompiler LATEX.EXE (LaTeX=>PS=>PDF)
%%%%%%%%%%%%%%%%%%%%%%%%%%%%%%%%%%%%%%%%%%%%%%%%%%%%%%%%%%%%%%%%%%%%%%%%
% Ueberarbeitung für pdflatex (LaTeX=>PDF)
%%%%%%%%%%%%%%%%%%%%%%%%%%%%%%%%%%%%%%%%%%%%%%%%%%%%%%%%%%%%%%%%%%%%%%%%
% Vorlage Changelog:
% 10.09.2015 (Matthias Liebeck): Nummerierung des Inhaltsverzeichnis nun römisch, Beispiel für einen Anhang eingebaut, \raggedbottom hinter sections eingefügt
% 11.07.2018 (Matthias Liebeck): Ersetzung des Bibliographiestils, Einsatz von Biber
% 04.09.2018 (Matthias Liebeck):
%   * Bibtex: unnötige Bibtexfelder beim Rendern ausblenden (thx @ Markus Brenneis)
%   * ngerman: "et al." im BibTeX für drei oder mehr Autoren
%   * Neuer Befehl \sectionforcestartright: Sections immer rechts beginnen (thx @ Philipp Grawe)
%   * ngerman: Deutsche Anführungszeichen im Literaturverzeichnis (thx @ Markus Brenneis)
%   * ngerman: Deutsche Anführungszeichen im Literaturverzeichnis (thx @ Markus Brenneis)
% 16.10.2018 (Matthias Liebeck): Zwei fixes an \sectionforcestartright (thx @ Markus Brenneis)
%%%%%%%%%%%%%%%%%%%%%%%%%%%%%%%%%%%%%%%%%%%%%%%%%%%%%%%%%%%%%%%%%%%%%%%%
%%%% BEGINN EINSTELLUNG FUER DIE ARBEIT. UNBEDINGT ERFORDERLICH! %%%%%%%
%%%%%%%%%%%%%%%%%%%%%%%%%%%%%%%%%%%%%%%%%%%%%%%%%%%%%%%%%%%%%%%%%%%%%%%%
% Geben Sie Ihren Namen hier an:
\newcommand{\bearbeiter}{Gaius Julius Caesar}

% Geben Sie hier den Titel Ihrer Arbeit an:
\newcommand{\titel}{De Bello Gallico}

% Geben Sie das Datum des Beginns und Ende der Bachelorarbeit ein:
\newcommand{\beginndatum}{01. April 2010}
\newcommand{\abgabedatum}{31.~September~2010}

% Geben Sie die Namen des Erst- und Zweitgutachters an:
\newcommand{\erstgutachter}{Prof. Dr.~Stefan Conrad}
\newcommand{\zweitgutachter}{Prof. Dr.~Stefan Conrad}

% Falls Sie die Arbeit zweiseitig ausdrucken wollen,
% benutzen Sie die folgende Zeile mit
% \AN fuer zweiseitigen Druck
% \AUS fuer einseitigen Druck
\newcommand{\zweiseitig}{\AN}

% Falls Sections immer rechts beginnen sollen. Gerade für Masterarbeiten
% interessant. Bei kurzen Bachelorarbeiten eher weniger zu verwenden.
\newcommand{\sectionforcestartright}{\AUS}
%\newcommand{\sectionforcestartright}{\AN}

% Falls die Arbeit in englischer Sprache verfasst 
% werden soll, dann benutzen Sie die folgende Zeile mit
% englisch fuer englische Sprache
% deutsch fuer deutsche Sprache
\newcommand{\sprache}{deutsch}

% Hier wird eingestellt, ob es sich bei der Arbeit um eine Bachelor- 
% oder Masterarbeit handelt (unpassendes auskommentieren!):
\newcommand{\arbeit}{Bachelorarbeit}
%~ \newcommand{\arbeit}{Masterarbeit}


%%%%%%%%%%%%%%%%%%%%%%%%%%%%%%%%%%%%%%%%%%%%%%%%%%%%%%%%%%%%%%%%%%%%%%%%
%%%% ENDE EINSTELLUNGEN %%%%%%%%%%%%%%%%%%%%%%%%%%%%%%%%%%%%%%%%%%%%%%%%
%%%%%%%%%%%%%%%%%%%%%%%%%%%%%%%%%%%%%%%%%%%%%%%%%%%%%%%%%%%%%%%%%%%%%%%%

% Die folgende Zeile NICHT EDITIEREN oder loeschen


%%%%%%%%%%%%%%%%%%%%%%%%%%%%%%%%%%%%%%%%%%%%%%%%%%%%%%%%%%%
% Obere Titelmakros. Editieren Sie diese Datei nur, wenn
% Sie sich ABSOLUT sicher sind, was Sie da tun!!!
% (Z.B. zum Abaendern der BA-Vorlage in eine MA-Vorlage)
% Uni Duesseldorf
% Lehrstuhl fuer Datenbanken und Informationssysteme
% Version 2.2 - 2.3.2010
%%%%%%%%%%%%%%%%%%%%%%%%%%%%%%%%%%%%%%%%%%%%%%%%%%%%%%%%%%%
\newcommand{\AN}{twoside}
\newcommand{\AUS}{}
%\newcommand{\englisch}{}
%\newcommand{\deutsch}{\usepackage[german]{babel}}

%% Die folgenden auskommentierten Optionen dienen der automatischen
%% Erkennung des Latex-Kompilers und dem Setzen der davon abhängigen
%% Einstellungen. Bei Problem z.B. mit dem Einbinden von verschiedenen
%% Grafiktypen bei Verwendung von PdfLatex oder Latex, einfach die
%% verschiedenen \usepackage(s) ausprobieren. (Mit diesen Einstellungen
%% funktionierte diese Vorlage bei der Verwenundg von latex.exe als
%% Kompiler bei den meisten Studierenden.)

%\newif\ifpdf \ifx\pdfoutput\undefined
%\pdffalse % we are not running pdflatex
%\else
%\pdfoutput=1 % we are running pdflatex
%\pdfcompresslevel=9 % compression level for text and image;
%\pdftrue \fi

\documentclass[11pt,a4paper, \zweiseitig]{article}



%\usepackage[iso]{umlaute}
\usepackage[utf8]{inputenc}
\usepackage{palatino} % palatino Schriftart
%\usepackage{makeidx} % um ein Index zu erstellen
\usepackage[nottoc]{tocbibind}
\usepackage[T1]{fontenc} %fuer richtige Trennung bei Umlauten
\usepackage{fancybox} % fuer die Rahmen
\usepackage{shortvrb}
\usepackage{url}
\usepackage{xcolor}
\usepackage[colorlinks,citecolor=blue,linkcolor=black]{hyperref} %anklickbares Inhaltsverzeichnis
\usepackage[style=authoryear,natbib=true,backend=biber,mincitenames=1,maxcitenames=2,maxbibnames=99,uniquelist=false,dashed=false]{biblatex}

%% Extra packages
\usepackage[colorinlistoftodos,obeyFinal]{todonotes}
\usepackage{graphicx}
\usepackage{booktabs}
\usepackage{multirow}

%% Math
\usepackage{amsmath}
\usepackage{amssymb}

\DeclareMathOperator*{\minimize}{minimize}

% https://tex.stackexchange.com/a/334703/8850
\AtEveryBibitem{%
  \clearfield{issn}
  \clearfield{isbn}
  \clearfield{doi}
  \clearfield{location}
  \clearlist{location}
  \clearlist{address}

  \ifentrytype{online}{}{% Remove url except for @online
    \clearfield{url}
  }
}

% Falls es bei \citet ein Komma zwischen Name und Jahr gibt:
% https://tex.stackexchange.com/questions/312539/unwanted-comma-between-author-and-year-using-citet-command
% (thx @ Markus Brenneis)
%\DeclareDelimFormat[cbx@textcite]{nameyeardelim}{\addspace}



\usepackage{ifthen}
\ifthenelse{\equal{\sprache}{deutsch}}{
 \usepackage[ngerman]{babel}
 % Bibtex u.a -> et al.
 \DefineBibliographyStrings{ngerman}{
   andothers = {{et\,al\adddot}},
 }
 \usepackage[autostyle, german=quotes]{csquotes} % Deutsche Anführungszeichen im Literaturverzeichnis (thx @ Markus Brenneis)
 }{}

\usepackage{a4wide} % ganze A4 Weite verwenden



%\ifpdf
%\usepackage[pdftex,xdvi]{graphicx}
%\usepackage{thumbpdf} %thumbs fuer Pdf
%\usepackage[pdfstartview=FitV]{hyperref} %anklickbares Inhaltsverzeichnis
%\else
%\usepackage[dvips,xdvi]{graphicx}
\usepackage{graphicx}

%\fi

\newcommand{\redt}[1] {
  \textcolor{red}{#1}}

\newcommand{\oranget}[1] {
  \textcolor{orange}{#1}}

\newcommand{\purplet}[1] {
  \textcolor{purple}{#1}}

%%%%%%%%%%%%%%%%%%%%%%% Massangaben fuer die Arbeit %%%%%%%%%%%%%%%
\setlength{\textwidth}{15cm}

\setlength{\oddsidemargin}{35mm}
\setlength{\evensidemargin}{25mm}

\addtolength{\oddsidemargin}{-1in}
\addtolength{\evensidemargin}{-1in}

\addbibresource{references.bib}

%\makeindex

\begin{document}

%\setcounter{secnumdepth}{4} %Nummerieren bis in die 4. Ebene
%\setcounter{tocdepth}{4} %Inhaltsverzeichnis bis zur 4. Ebene

\pagestyle{headings}

\sloppy % LaTeX ist dann nicht so streng mit der Silbentrennung
%~ \MakeShortVerb{\§}

\parindent0mm
\parskip0.5em


{
\textwidth170mm 
\oddsidemargin30mm 
\evensidemargin30mm 
\addtolength{\oddsidemargin}{-1in}
\addtolength{\evensidemargin}{-1in}

\parskip0pt plus2pt

% Die Raender muessen eventuell fuer jeden Drucker individuell eingestellt
% werden. Dazu sind die Werte fuer die Abstaende `\oben' und `\links' zu
% aendern, die von mir auf jeweils 0mm eingestellt wurden.

%\newlength{\links} \setlength{\links}{10mm}  % hier abzuaendern
%\addtolength{\oddsidemargin}{\links}
%\addtolength{\evensidemargin}{\links}

\begin{titlepage}
\vspace*{-1.5cm}
  \raisebox{17mm}{
    \begin{minipage}[t]{70mm}
      \begin{center}
        %\selectlanguage{german}
        {\Large INSTITUT FÜR INFORMATIK\\}
        {\normalsize
          Machine Learning\\
        }
        \vspace{3mm}
        {\small Universitätsstr. 1 \hspace{5ex} D--40225 Düsseldorf\\}
     \end{center}
    \end{minipage}
  }
  \hfill
  \includegraphics[width=130pt]{bilder/HHU_Logo}
  \vspace{14em}

% Titel
  \begin{center}
      	\baselineskip=55pt
    	\textbf{\huge \titel}
  	 	\baselineskip=0 pt
   \end{center}

  %\vspace{7em}

\vfill

% Autor
  \begin{center}
    \textbf{\Large
      \bearbeiter
    }
  \end{center}

  \vspace{35mm}
 
% Prüfungsordnungs-Angaben
  \begin{center}
    %\selectlanguage{german}
    
%%%%%%%%%%%%%%%%%%%%%%%%%%%%%%%%%%%%%%%%%%%%%%%%%%%%%%%%%%%%%%%%%%%%%%%%%
% Ja, richtig, hier kann die BA-Vorlage zur MA-Vorlage gemacht werden...
% (nicht mehr nötig!)
%%%%%%%%%%%%%%%%%%%%%%%%%%%%%%%%%%%%%%%%%%%%%%%%%%%%%%%%%%%%%%%%%%%%%%%%%
    {\Large \arbeit}

    \vspace{2em}

    \begin{tabular}[t]{ll}
      Beginn der Arbeit:& \beginndatum \\
      Abgabe der Arbeit:& \abgabedatum \\
      Gutachter:         & \erstgutachter \\
                         & \zweitgutachter \\
    \end{tabular}
  \end{center}

\end{titlepage}

}

%%%%%%%%%%%%%%%%%%%%%%%%%%%%%%%%%%%%%%%%%%%%%%%%%%%%%%%%%%%%%%%%%%%%%
\clearpage
\begin{titlepage}
  ~                % eine leere Seite hinter dem Deckblatt
\end{titlepage}
%%%%%%%%%%%%%%%%%%%%%%%%%%%%%%%%%%%%%%%%%%%%%%%%%%%%%%%%%%%%%%%%%%%%%
\clearpage
\begin{titlepage}
\vspace*{\fill}

\section*{Erklärung}

%%%%%%%%%%%%%%%%%%%%%%%%%%%%%%%%%%%%%%%%%%%%%%%%%%%%%%%%%%%
% Und hier ebenfalls ggf. BA durch MA ersetzen...
% (Auch nicht mehr nötig!)
%%%%%%%%%%%%%%%%%%%%%%%%%%%%%%%%%%%%%%%%%%%%%%%%%%%%%%%%%%%

Hiermit versichere ich, dass ich diese \arbeit{}
selbstständig verfasst habe. Ich habe dazu keine anderen als die
angegebenen Quellen und Hilfsmittel verwendet.

\vspace{25 mm}

\begin{tabular}{lc}
Düsseldorf, den \abgabedatum \hspace*{2cm} & \underline{\hspace{6cm}}\\
& \bearbeiter
\end{tabular}

\vspace*{\fill}
\end{titlepage}

%%%%%%%%%%%%%%%%%%%%%%%%%%%%%%%%%%%%%%%%%%%%%%%%%%%%%%%%%%%%%%%%%%%%%
% Leerseite bei zweiseitigem Druck
%%%%%%%%%%%%%%%%%%%%%%%%%%%%%%%%%%%%%%%%%%%%%%%%%%%%%%%%%%%%%%%%%%%%%

\ifthenelse{\equal{\zweiseitig}{twoside}}{\clearpage\begin{titlepage}
~\end{titlepage}}{}

%%%%%%%%%%%%%%%%%%%%%%%%%%%%%%%%%%%%%%%%%%%%%%%%%%%%%%%%%%%%%%%%%%%%%
\clearpage
\begin{titlepage}

%%% Die folgende Zeile nicht ändern!
\section*{\ifthenelse{\equal{\sprache}{deutsch}}{Zusammenfassung}{Abstract}}
%%% Zusammenfassung:

This bachelor thesis extensively evaluates the vulnerability of capsule network to various adversarial attacks.
Capsule networks are a form of neural network in which neurons are grouped in capsules whose activity vectors describe different properties of the same entity.
Recent work attributes them more robustness to adversarial attacks than other types of neural networks.

To investigate those conjectures we train four pairs of capsule networks and convolutional neural networks discriminatively on different image datasets.
Using multiple strong and diverse adversarial attack we compute adversarial examples for each combination of dataset and architecture.

By comparing the euclidean norms of the generated perturbations we find that capsule networks are vulnerable to adversarial attacks in a similar level as convolutional networks.
We discover in particular that capsule networks show no robustness to sufficiently strong white-box attacks.
Additionally we demonstrate limited transferability of the adversarial examples between the two types of architectures.

We study furthermore the structure of the adversarial perturbations and can observe that they seem to lie in a relatively low dimensional linear subspace in both cases.
Despite those similarities, using dimensionality reduction we detect inherent differences between the perturbations constructed for the different network types.

Finally we aim to identify the source of adversarial examples in the networks we trained.
Since capsule networks and convolutional partially use the same type of convolutional layer we focus on the disturbance on the activation of individual layers caused by adversarial attacks.
We find that susceptibility to attacks is not just ascribable to the first layers in a network but a result of deeper layers as well.
Specifically we find that the dynamic routing algorithm used in capsule networks does not establish capsule activation invariants to counteract adversarial attacks as previously conjectured.


%%%%%%%%%%%%%%%%%%%%%%%%%%%%%%%%%%%%%%%%%%%%%%%%
% Untere Titelmakros. Editieren Sie diese Datei nur, wenn Sie sich
% ABSOLUT sicher sind, was Sie da tun!!!
%%%%%%%%%%%%%%%%%%%%%%%%%%%%%%%%%%%%%%%%%%%%%%%
\vspace*{\fill}
\end{titlepage}

%%%%%%%%%%%%%%%%%%%%%%%%%%%%%%%%%%%%%%%%%%%%%%%%%%%%%%%%%%%%%%%%%%%%%
% Leerseite bei zweiseitigem Druck
%%%%%%%%%%%%%%%%%%%%%%%%%%%%%%%%%%%%%%%%%%%%%%%%%%%%%%%%%%%%%%%%%%%%%
\ifthenelse{\equal{\zweiseitig}{twoside}}
  {\clearpage\begin{titlepage}~\end{titlepage}}{}
%%%%%%%%%%%%%%%%%%%%%%%%%%%%%%%%%%%%%%%%%%%%%%%%%%%%%%%%%%%%%%%%%%%%%
\clearpage \setcounter{page}{1}
\pagenumbering{roman}
\setcounter{tocdepth}{2}
\tableofcontents

%\enlargethispage{\baselineskip}
\clearpage
%%%%%%%%%%%%%%%%%%%%%%%%%%%%%%%%%%%%%%%%%%%%%%%%%%%%%%%%%%%%%%%%%%%%%
% Leere Seite, falls Inhaltsverzeichnis mit ungerader Seitenzahl und 
% doppelseitiger Druck
%%%%%%%%%%%%%%%%%%%%%%%%%%%%%%%%%%%%%%%%%%%%%%%%%%%%%%%%%%%%%%%%%%%%%
\ifthenelse{ \( \equal{\zweiseitig}{twoside} \and \not \isodd{\value{page}} \)}
	{\pagebreak \thispagestyle{empty} \cleardoublepage}{\clearpage}


% Kapitel soll bei doppelseitigem Druck immer auf der rechten (ungeraden) Seite anfangen (thx @ Philipp Grawe)
% https://tex.stackexchange.com/a/223387
\ifthenelse{\( \equal{\sectionforcestartright}{twoside} \)}
 {\let\oldsection\section % Store \section in \oldsection
   \renewcommand{\section}{\cleardoublepage\oldsection}}
 {}
\pagenumbering{arabic}
\setcounter{page}{1}

%%%%%%%%%%%%%%%%%%%%%%%%%%%%%%%%%%%%%%%%%%%%%%%%%%%%%%%%%%%%%%%%%%%%%%%%
%%%% BEGINN TEXTTEIL %%%%%%%%%%%%%%%%%%%%%%%%%%%%%%%%%%%%%%%%%%%%%%%%%%%
%%%%%%%%%%%%%%%%%%%%%%%%%%%%%%%%%%%%%%%%%%%%%%%%%%%%%%%%%%%%%%%%%%%%%%%%

%%%%%%%%%%%%%%%%%%%%%%%%%%%%%%%%%%%%%%%%%%%%%%%%%%%%%%%%%%%%%%%%%%%%%%%%
% Text entweder direkt hier hinein schreiben oder, im Sinne der
% besseren Uebersichtlich- und Bearbeitbarkeit mittels \input die
% einzelnen Textteile hier einbinden.
%%%%%%%%%%%%%%%%%%%%%%%%%%%%%%%%%%%%%%%%%%%%%%%%%%%%%%%%%%%%%%%%%%%%%%%%

\section{De Bello Gallico}\raggedbottom 

\subsection{Gallia est omnis divisa} Gallia est omnis divisa in partes tres,
quarum unam
incolunt Belgae, aliam Aquitani, tertiam qui ipsorum lingua
Celtae, nostra Galli appellantur. Hi omnes lingua, institutis,
legibus inter se differunt. Gallos ab Aquitanis Garumna flumen, a
Belgis Matrona et Sequana dividit. Horum omnium fortissimi sunt
Belgae, propterea quod a cultu atque humanitate provinciae
longissime absunt, minimeque ad eos mercatores saepe commeant
atque ea quae ad effeminandos animos pertinent important,
proximique sunt Germanis, qui trans Rhenum incolunt, quibuscum
continenter bellum gerunt. Qua de causa Helvetii quoque reliquos
Gallos virtute praecedunt, quod fere cotidianis proeliis cum
Germanis contendunt, cum aut suis finibus eos prohibent aut ipsi
in eorum finibus bellum gerunt. Eorum una, pars, quam Gallos
obtinere dictum est, initium capit a flumine Rhodano, continetur
Garumna flumine, Oceano, finibus Belgarum, attingit etiam ab
Sequanis et Helvetiis flumen Rhenum, vergit ad septentriones.
Belgae ab extremis Galliae finibus oriuntur, pertinent ad
inferiorem partem fluminis Rheni, spectant in septentrionem et
orientem solem. Aquitania a Garumna flumine ad Pyrenaeos montes et
eam partem Oceani quae est ad Hispaniam pertinet; spectat inter
occasum solis et septentriones.


\subsection{Apud Helvetios longe} Apud Helvetios longe nobilissimus fuit et
ditissimus
Orgetorix. Is M. Messala, [et P.] M. Pisone consulibus regni
cupiditate inductus coniurationem nobilitatis fecit et civitati
persuasit ut de finibus suis cum omnibus copiis exirent: perfacile
esse, cum virtute omnibus praestarent, totius Galliae imperio
potiri. Id hoc facilius iis persuasit, quod undique loci natura
Helvetii continentur: una ex parte flumine Rheno latissimo atque
altissimo, qui agrum Helvetium a Germanis dividit; altera ex parte
monte Iura altissimo, qui est inter Sequanos et Helvetios; tertia
lacu Lemanno et flumine Rhodano, qui provinciam nostram ab
Helvetiis dividit. His rebus fiebat ut et minus late vagarentur et
minus facile finitimis bellum inferre possent; qua ex parte
homines bellandi cupidi magno dolore adficiebantur. Pro
multitudine autem hominum et pro gloria belli atque fortitudinis
angustos se fines habere arbitrabantur, qui in longitudinem milia
passuum CCXL, in latitudinem CLXXX patebant.

His rebus adducti et auctoritate Orgetorigis permoti constituerunt
ea quae ad proficiscendum pertinerent comparare, iumentorum et
carrorum quam maximum numerum coemere, sementes quam maximas
facere, ut in itinere copia frumenti suppeteret, cum proximis
civitatibus pacem et amicitiam confirmare. Ad eas res conficiendas
biennium sibi satis esse duxerunt; in tertium annum profectionem
lege confirmant. Ad eas res conficiendas Orgetorix deligitur. Is
sibi legationem ad civitates suscipit. In eo itinere persuadet
Castico, Catamantaloedis filio, Sequano, cuius pater regnum in
Sequanis multos annos obtinuerat et a senatu populi Romani amicus
appellatus erat, ut regnum in civitate sua occuparet, quod pater
ante habuerit; itemque Dumnorigi Haeduo, fratri Diviciaci, qui eo
tempore principatum in civitate obtinebat ac maxime plebi acceptus
erat, ut idem conaretur persuadet eique filiam suam in matrimonium
dat. Perfacile factu esse illis probat conata perficere, propterea
quod ipse suae civitatis imperium obtenturus esset: non esse
dubium quin totius Galliae plurimum Helvetii possent; se suis
copiis suoque exercitu illis regna conciliaturum confirmat. Hac
oratione adducti inter se fidem et ius iurandum dant et regno
occupato per tres potentissimos ac firmissimos populos totius
Galliae sese potiri posse sperant.

\citet{Con97} hat ein Buch geschrieben. Es gibt auch andere Arbeiten \citep{PeHe97} die referenziert sind. In Abbildung \ref{fig_Gallien} ist ein Sachverhalt dargestellt.


1 Autor: \citet{Con97} \hspace*{1cm} \citep{Con97}\\
2 Autoren: \citet{IWNLP} \hspace*{1cm} \citep{IWNLP}\\
3 Autoren: \citet{liebeck-esau-conrad:2016:ArgMining2016} \hspace*{1cm} \citep{liebeck-esau-conrad:2016:ArgMining2016}

Online resource: \citet{ILSVRC2016}

\textbf{quotes}:\\
\glqq quote\grqq.

\begin{figure}[htb]
\begin{center}
  \includegraphics[width=175pt, angle=270]{bilder/Galia}
  \caption{Gallien zur Zeit Caesars}\label{fig_Gallien}
\end{center}
\end{figure}


\begin{table}[htb]
\begin{center}
\begin{tabular}{|l|l|l|}
\hline
Jahr &  Erster Consul & Zweiter Consul\\
\hline \hline
1 & C. Caesar         & L. Aemilius Paullus\\
2 & P. Vinicius       & P. Alfenus Varus\\
3 & L. Aelius Lamia   & M. Servilius\\
4 & Sex. Aelius Catus &  C. Sentius Saturninus\\
5 & L. Valerius Messalla& Cn. Cornelius Cinna \\
suff. & C. Vibius Postumus &  C. Ateius Capito\\
6 & M. Aemilius Lepidus & L. Arruntius\\
\hline
\end{tabular}
 \caption{Römische Konsulen}\label{tab_Konsulen}
\end{center}
\end{table}


\pagebreak

\section{De Bello Hispaniensi}\raggedbottom 

Pharnace superato, Africa recepta, qui ex his proeliis cum
adulescente Cn. Pompeio profugissent, cum . . . et ulterioris
Hispaniae potitus esset, dum Caesar muneribus dandis in Italia
detinetur, . . . quo facilius praesidia contra compararet,
Pompeius in fidem uniuscuiusque civitatis confugere coepit. Ita
partim precibus partim vi bene magna comparata manu provinciam
vastare. Quibus in rebus non nullae civitates sua sponte auxilia
mittebant, item non nullae portas contra cludebant. Ex quibus si
qua oppida vi ceperat, cum aliquis ex ea civitate optime de Cn.
Pompeio meritus civis esset, propter pecuniae magnitudinem alia
qua ei inferebatur causa, ut eo de medio sublato ex eius pecunia
latronum largitio fieret. Ita pacis commoda hoste +hortato+
maiores augebantur copiae. +Hoc crebris nuntiis in Italiam missis
civitates contrariae Pompeio+ auxilia sibi depostulabant.

C. Caesar dictator tertio, designatus dictator quarto multis
+iterante diebus coniectis+ cum celeri festinatione ad bellum
conficiendum in Hispaniam cum venisset, legatique Cordubenses, qui
a Cn. Pompelo discessissent, Caesari obviam venissent, a quibus
nuntiabatur nocturno tempore oppidum Cordubam capi posse, quod nec
opinantibus adversariis eius provinciae potitus esset, simulque
quod tabellariis, qui a Cn. Pompeio dispositi omnibus locis
essent, qui certiorem Cn. Pompeium de Caesaris adventu facerent .
. . multa praeterea veri similia proponebant. Quibus rebus
adductus quos legatos ante exercitui praefecerat Q. Pedium et Q.
Fabium Maximum de suo adventu facit certiores, utque sibi
equitatus qui ex provincia fuisset praesidio esset. Ad quos
celerius quam ipsi opinati sunt appropinquavit neque, ut ipse
voluit, equitatum sibi praesidio habuit.

Erat idem temporis Sex. Pompeius frater qui cum praesidio Cordubam
tenebat, quod eius provinciae caput esse existimabatur; ipse autem
Cn. Pompeius adulescens Uliam oppidum oppugnabat et fere iam
aliquot mensibus ibi detinebatur. Quo ex oppido cognito Caesaris
adventu legati clam praesidia Cn. Pompei Caesarem cum adissent,
petere coeperunt uti sibi primo quoque tempore subsidium mitteret.
Caesar - eam civitatem omni tempore optime de populo Romano
meritam esse - celeriter sex cohortis secunda vigilia iubet
proficisci, pari equites numero. Quibus praefecit hominem eius
provinciae notum et non parum scientem, L. Vibiurn Paciaecum. Qui
cum ad Cn. P praesidia venisset, incidit idem temporis ut
tempestate adversa vehementique vento adflictaretur; aditusque vis
tempestatis ita obscurabat ut vix proximum agnoscere possent.
Cuius incommodum summam utilitatem ipsis praebebat. Ita cum ad eum
locum venerunt, iubet binos equites conscendere, et recta per
adversariorum praesidia ad oppidum contendunt. Mediisque eorum
praesidiis cum essent, cum quaereretur qui essent unus ex nostris
respondit, ut sileat verbum facere: nam id temporis conari ad
murum accedere, ut oppidum capiant; et partim tempestate impediti
vigiles non poterant diligentiam praestare, partim illo responso
deterrebantur. Cum ad portam appropinquassent, signo dato ab
oppidanis sunt reccepti, et pedites dispositi partim ibi
remanserunt, equites clamore facto eruptionem in adversariorum
castra fecerunt.

\pagebreak
\section{Weiteres Kapitel}\raggedbottom 
\subsection{Unterkapitel}
\subsection{Unterkapitel}

%%%%%%%%%%%%%%%%%%%%%%%%%%%%%%%%%%%%%%%%%%%%%%%%%%%%%%%%%%%%%%%%%%%%%%%%
%%%% ENDE TEXTTEIL %%%%%%%%%%%%%%%%%%%%%%%%%%%%%%%%%%%%%%%%%%%%%%%%%%%%%
%%%%%%%%%%%%%%%%%%%%%%%%%%%%%%%%%%%%%%%%%%%%%%%%%%%%%%%%%%%%%%%%%%%%%%%%

\clearpage

% Entfernen Sie das Kommentar aus der nachfolgenden Zeile, falls Sie einen Anhang in der Arbeit verwenden wollen. Beachten Sie, dass Sie sich im Verlauf der Arbeit mit \ref{...} (z.B. \ref{anhang:zusatz1}) auf den Anhang beziehen.
%\newpage
\appendix
\section{Anhang}

\subsection*{Zusatzteil 1} \label{anhang:zusatz1}

Dies ist ein Anhang.

\clearpage

%\bibliography{references}
%\bibliographystyle{alphadin}
\DeclareNameAlias{sortname}{first-last}
\printbibliography[heading=bibintoc, title=References]
%\vspace*{\fill}

\clearpage

\listoffigures

\listoftables

%\pagebreak

%\printindex
\end{document}
